\chapter{Discussion}
\label{chap:discussion}

This chapter interprets the two observational causal-analysis testbeds characterized in Chapter~\ref{chap:results} within the study design and assumptions established in Chapters~\ref{chap:problem-definition-study-design}--\ref{chap:experimental-design}.  The central distinction is between evidence that characterizes a fixed observational resource and evidence that would score an estimator against known causal truth.  The latter is unavailable here.  Prediction results, descriptive matched-pair comparisons, model-estimated contrasts, clinical reference points, and diagnostics are therefore kept non-equivalent and bounded by their sources.

\section{Answering the Research Questions}
\label{sec:discussion:answer-research-questions}
% Section ID: C11.1.
% Evidence map: see thesis-writing/planning/chapter_evidence_map.md.

\subsection{Main Research Question}
\label{subsec:discussion:main-rq}

The main research question asked: \emph{How can clinically grounded observational causal-analysis testbeds be constructed from source-observed irregular ICU records---without simulating patient records, treatment or exposure-assignment mechanisms, or outcomes---and what converging evidence can characterize their usefulness and limitations for evaluating causal estimators?}

CliniCause constructed two resources separately from MIMIC-III and PhysioNet 2012.  Source-recorded measurements and mortality pass through explicit preprocessing and identifier contracts, while project-selected proxies, rules, cohorts, DAGs, adjustment assumptions, estimator interfaces, and provenance form an explicit representation layer.  A schema-bounded LLM-assisted protocol operated only at design time; it did not access patient records, assign exposures, simulate outcomes, or estimate effects.  The integrated interfaces make the resources inspectable.

Because neither resource contains a counterfactual answer key, converging, non-equivalent evidence characterizes rather than scores it.  Recoverability, estimator agreement, clinical comparison, support, disruption, outcome-downsampling, sensitivity, and provenance show persistence and failure modes under a fixed representation.  Agreement is not accuracy; the resources support retrospective evaluation and hypothesis generation, not clinical validation, treatment recommendation, or deployment.

\subsection{SRQ-1---Representation Instantiation}
\label{subsec:discussion:data-contracts-integration}

The representation designs began with dataset schemas, the bounded literature corpus, and structured design-time LLM assistance.  Project selection converted candidate ontologies, rule families, missingness considerations, and DAGs into deterministic source.  The final rule code and stable \texttt{LAT\_*} fields make the implemented proxy states inspectable and replaceable; they are not diagnoses or discovered biological variables.  No patient record was supplied to the LLM, and the execution path used no runtime LLM call.

Instantiation remained dataset specific.  MIMIC-III and PhysioNet expose different variables, frequencies, missingness processes, ontologies, thresholds, and graph structures, so similarly named proxy states are not assumed to be equivalent.  The available construct evidence is plausibility only: no chart adjudication, blinded multi-clinician rating, inter-rater analysis, clinician-only comparison, or clinical validation was performed.  Separate construction and the no-pooling policy preserve those differences instead of converting nominally similar labels into one common exposure.

\subsection{SRQ-2---Contracts, Aggregation, DAG, and Provenance}
\label{subsec:discussion:proxy-prediction-aggregation}

The causal repository uses \texttt{[ts, oc, ts\_ids]}, while the predictive workflow uses a split-aware artifact with explicit train, validation, and test identifiers.  Prediction exports add paired probability and binary fields followed by binary-only voter files.  Boundary checks prevent silent identifier, split, schema, and cohort conflation, although processed-pickle hashes, producing commands, and final split lineage remain incomplete.

Majority voting aligns identifiers, checks identical binary schemas, and applies the implemented tie rule to one rule-derived source and four thresholded predictions---GRU, GRU-D, STraTS, and TCN.  This fixed algorithmic aggregate is not clinical consensus or a learned source-quality model; shared rule error can persist because all learned voters target the same rules and the rule source itself participates \cite{ratner_et_al_2016_data_programming}.

Dataset-specific source-coded DAGs define exposure-specific observed adjustment sets from available columns.  They make the project's assumptions inspectable but were neither learned nor clinically validated, and a computed set does not prove exchangeability or the absence of unmeasured confounding \cite{pearl_1995_causal_diagrams}.  Source paths, status classes, manifests, and checked tables expose the lineage and diagnostic gaps.  These contracts, aggregation rules, graph interfaces, and provenance records are enabling mechanisms for testbed integrity; they are not the primary intellectual object and do not establish scientific validity by themselves.

\subsection{SRQ-3---Recoverability and Mortality-Relevant Information}
\label{subsec:discussion:recoverability-mortality}

The four-model comparison shows that rule-derived proxy labels are recoverable from irregular source records, with archived AUROC spanning 0.867--0.918.  STraTS led MIMIC and GRU-D led PhysioNet; changing lower-order rankings preclude a universal ordering.  Recoverability depends on the measurements, missingness, labels, and representation, and the local tasks do not reproduce the published benchmark tasks exactly \cite{sun_2026_review_irregular_medical_timeseries,tipirneni2022strats,che2018grud}.

The admitted MIMIC-only proxy-vector analysis adds mortality-relevant predictive information: logistic regression and the MLP reached test AUROC 0.826 and 0.831, respectively, for source-observed in-hospital mortality.  This is prognostic association, not causal evidence.  The analogous PhysioNet number remains withheld because checked-thesis and paper voter/cohort lineages conflict.  That asymmetry is a provenance limitation to preserve, not a result to reconcile rhetorically.  No uncertainty intervals, paired between-model tests, or clinical validation of the targets were available.

\subsection{SRQ-4---Converging Evidence}
\label{subsec:discussion:adjusted-effects-agreement}

\paragraph{Estimator triangulation and diagnostics.}
The hierarchy uses original cohorts as primary, CausalForestDML as the primary estimator, LinearDML as secondary, CausalPFN as exploratory, and outcome-downsampled analyses as robustness evidence.  The implemented representation and adjustment procedure produced positive primary forest means for all nine MIMIC exposures and nine of ten PhysioNet exposures, with shock negative.  These are construct- and resource-dependent model-estimated contrasts, not evidence that a proxy caused death, that changing it would change mortality, or that a ranking identifies intervention priorities.

CausalForestDML and LinearDML agreed in direction for all 19 original-cohort dataset--exposure pairs.  Across all three estimators, agreement was 18 of 19.  Pairwise within-resource rank correlations were 0.983--1.000 in MIMIC and 0.794--0.964 in PhysioNet, while pairwise RMSE was 1.35--2.57 and 1.08--2.04 percentage points, respectively.  These are distinct views of stability, not measures of estimator accuracy.  They show that most patterns were not unique to one fitted estimator under a fixed representation and population; they do not validate the DAG, establish equal estimands or uncertainty, or rule out residual confounding.

Within this pipeline, CausalPFN was a useful complementary estimator because it reproduced the prevailing direction across nearly all prespecified comparisons despite using a different modeling approach.  This result is bounded: magnitudes were not identical to DML magnitudes, and its archived sensitivity and permutation stages were intentionally absent.  CausalPFN therefore remains exploratory rather than an estimator with the same diagnostic support as the DML analyses.

PhysioNet shock is the central exception rather than a detail to be averaged away.  CausalForestDML and LinearDML were negative, CausalPFN was slightly positive, and the descriptive matching summary was positive.  The external shock contrast is not commensurate in construct, population, or horizon.  This combination warns about representation, support, estimator, and population dependence.  No estimator vote resolves it, and no protective clinical effect should be inferred from the negative DML estimates.

Matching provides a descriptive baseline and empirical-support warnings.  Fourteen of 15 successful original-cohort summaries shared the primary forest direction; failures, insufficient-pair flags, and varying Hamming distances remain visible.  Unchanged matching availability across original and downsampled populations documents that availability pattern but proves neither positivity nor balance.  Failure is missing support, not a zero estimate, and a matched-pair difference is not automatically an ATT or independent causal confirmation.

Sensitivity and permutation evidence adds layers rather than closure.  At the strict $\lvert z\rvert>2$ disruption threshold, all 36 MIMIC DML comparisons and 34 of 40 PhysioNet DML comparisons passed; these are sanity checks across exposure, DML estimator, and treatment/outcome permutation, not p-values or identification tests.  DML sensitivity outputs were saved-training direct, reconstructed, partial, or failed depending on the exposure, and reconstructed diagnostics are not estimator-native calculations.  The evidence classes therefore support the qualitative conclusion that broad agreement can coexist with omitted-confounding fragility, but not one unified numerical threshold or exposure count.  CausalPFN diagnostics were intentionally skipped.  Dedicated propensity, common-support, and balance plots were not archived; positivity remains unestablished \cite{crump_et_al_2009_limited_overlap,cinelli_hazlett_2020_sensitivity,chernozhukov_et_al_2026_ovb_causal_ml}.  Agreement and fragility can coexist, and sensitivity prevents estimator convergence from being treated as closure.

\subsection{Cross-Dataset Interpretation and Robustness}
\label{subsec:discussion:cross-dataset-robustness}

The original causal-analysis populations contained 26,845 MIMIC records and 7,993 PhysioNet records, so the MIMIC analysis population was approximately 3.36 times larger.  A larger analysis population can provide more information for nuisance-model fitting and modeled effect heterogeneity, and it may contribute to numerical stability.  It does not demonstrate better data quality, make the MIMIC findings more valid, or erase differences in era, variables, measurement systems, missingness, proxy rules, DAGs, or patient composition.  Predictive training and causal analysis also use distinct artifact contracts, so their populations should not be assumed identical merely because some row counts align.

Operation of the same representation, prediction, aggregation, graph-guided adjustment, and estimator sequence on both datasets demonstrates workflow portability.  It does not make similarly named states equivalent.  Across nine approximate mappings including shock, three-estimator-average ranks had moderate correspondence ($\rho=0.533$) and eight mappings shared direction.  Dataset-specific rules and graphs and MIMIC's combined domains make the no-pooling policy essential; cross-resource differences characterize contextual dependence, not pooled effects or biological inconsistency.

Outcome-based downsampling retained outcome-positive records and sampled outcome-negative records, changing both population and prevalence.  Direction was preserved in 55 of 57 matched comparisons, documenting broad sign stability under this specified perturbation, but every paired mean changed and two PhysioNet comparisons changed sign.  The procedure is not transport or reweighting; the means need not share a population target, and stability does not validate identification.

The clinical-literature comparison must be interpreted symmetrically.  Renal/AKI, hepatic/bilirubin, and cardiac injury were broadly compatible or overlapping, while inflammation/sepsis, shock, coagulation/hematologic, respiratory, neurologic, and metabolic were discrepant; global severity had no defensible numerical comparator.  External definitions, severity, populations, adjustments, and horizons differ.  Agreement is contextual corroboration, not confirmation, and no external value is commensurate causal ground truth for a CliniCause proxy.

\subsection{Relationship to Prior Work and Methodological Implications}
\label{subsec:discussion:prior-work-implications}

The primary contribution is the construction and characterization of two observational causal-analysis testbeds.  Integrated representation, prediction, aggregation, DAG adjustment, estimator, diagnostic, and provenance interfaces make them inspectable by preventing silent identifier, split, cohort, and construct conflation and retaining both agreement and failure modes.

CliniCause complements controlled benchmarks.  Answer-key resources score error under specified mechanisms; observational testbeds retain source-recorded measurements, missingness patterns, care-process traces, covariates, and outcomes while exposing representation-dependent support, estimator, and provenance behavior without measuring causal error against truth.

This framing follows target-trial and well-defined-intervention principles \cite{hernan_robins_2016_target_trial,hernan_taubman_2008_well_defined_interventions} and ICU guidance on question definition, confounding, data preparation, and reporting \cite{smit_2023_causal_inference_icu_scoping_review}.  These sources provide context, not validation of local labels, graphs, estimates, or execution; CliniCause remains an auditable research workflow rather than a validated clinical methodology.

\subsection{SRQ-5---Limitations and Appropriate Use}
\label{subsec:discussion:clinical-meaning}

The proxy states summarize clinically inspired patterns and are analytically useful for organizing retrospective evidence.  They may support hypothesis generation and help identify proxy-state contrasts that merit targeted clinical and causal investigation.  High prediction performance against rule labels does not establish clinical accuracy, however, and a positive adjusted proxy-state contrast does not imply that clinicians should attempt to manipulate the proxy directly.  The system did not undergo prospective validation, external clinical validation, or clinical impact analysis.  Its causal quantities are not ready for bedside decision support, and no patient-level treatment recommendation follows from the results.

\paragraph{Principal limitations.}
The resources have no known causal truth and are \emph{empirically supported with substantial unresolved boundaries}.  The canonical construct, identification, support, generalizability, provenance, LLM-design, ethics, and deployment limitations follow.

\section{Limitations and Threats to Validity}
\label{sec:discussion:limitations-threats-validity}
% Section ID: C11.2.
% Evidence map: see thesis-writing/planning/chapter_evidence_map.md.

\subsection{Construct and Measurement Validity}
\label{subsec:discussion:construct-measurement-validity}

The principal construct limitation is that clinically inspired rules approximate conditions but were not evaluated against chart-adjudicated diagnoses.  Available evidence consists of active rule code, binary artifact schemas, and limited descriptive validation summaries.  This makes the constructs transparent but does not establish their sensitivity, specificity, or equivalence across datasets.  Prediction can add classification error, and the mixed-source majority vote can preserve shared systematic error among models trained on the same rules.  Existing mitigations are deterministic definitions, source-code traceability, dataset-specific naming, and explicit proxy terminology; the residual risk remains high because clinical review and reference-standard validation are absent.

Binary exposure states also compress severity, duration, and timing.  A positive tag may combine clinically distinct trajectories, while a negative tag can reflect missing measurement rather than absence of a condition.  Proxy onset relative to covariates and mortality is not fully established, so exposure misclassification and temporal ambiguity can bias both prediction and adjusted analyses.  The processed outcome, \texttt{in\_hospital\_mortality}, is consistently required by the software, but its producing-artifact provenance is incomplete, and no competing-risk or post-discharge mortality analysis was performed.

The observation process is itself informative.  Measurement frequency may reflect clinical concern, resource use, or care pathways; models that encode masks or time gaps can exploit this signal without removing its causal implications.  Preprocessing and missingness-aware representation are therefore not equivalent to adjustment for the decision to measure.  This concern can affect proxy construction, predictive comparisons, and effect estimation simultaneously.

\subsection{Causal Identification and Estimand Validity}
\label{subsec:discussion:identification-estimand-validity}

These are the central constraints on causal interpretation.  Illness-state proxies are not necessarily well-defined interventions, so consistency and the meaning of contrasting exposure states may be difficult to specify.  No target-trial emulation was performed, and the pipeline does not implement a method for time-varying treatment--confounder feedback.  Exchangeability depends on a project-specified DAG, correct temporal ordering, and measured variables; topology alone cannot establish patient-record timing.  Some archived observed adjustment sets are empty, every checked identifiability field remains unresolved, and unmeasured or misspecified confounding remains possible.

Positivity was assumed rather than established.  Matching availability, failures, pair counts, and distances provide indirect evidence, but the archive lacks dedicated propensity distributions, common-support plots, and pre/post-match balance summaries.  Interference is also assumed, not tested, and repeated admissions or shared ICU processes may complicate unit-level independence.

Estimand language remains deliberately restricted.  The mean model-estimated CATE over the analyzed sample is an aggregate of model outputs, not automatically a population ATE.  The descriptive matched-pair outcome difference is not automatically an ATT.  Normalized CATE was omitted because scaling by the sample outcome rate is unstable across populations and does not create a risk ratio or another standard clinical estimand.  These boundaries mitigate overstatement but do not resolve the underlying intervention-definition and identification problems.

\subsection{Statistical, Modeling, and Support Limitations}
\label{subsec:discussion:statistical-modeling-support}

Many displayed main estimates lack uncertainty intervals, no formal multiple-comparison strategy was documented, and this thesis did not add significance tests.  Absence of an interval cannot be interpreted as either significance or insignificance.  Model selection, nuisance-model adequacy, regularization, hyperparameters, and cross-fitting behavior remain assumptions whose impact was not exhaustively tested.  CausalForestDML, LinearDML, and CausalPFN do not provide interchangeable uncertainty, and directional agreement can coexist with material magnitude and ordering differences.

Support evidence is incomplete and representation dependent.  Some matching runs failed, some successful matches required larger Hamming distances, and some carried insufficient-pair warnings.  Sensitivity outputs were partial or failed for several exposure families; reconstructed or fallback diagnostics have different status from estimator-native outputs.  Permutation aggregates are limited sanity checks rather than formal inferential tests.  CausalPFN lacks the corresponding archived downstream diagnostic family.  The multi-estimator and diagnostic hierarchy exposes these risks, but it cannot remove them.

\subsection{External Validity and Generalizability}
\label{subsec:discussion:external-validity}

MIMIC-III and PhysioNet differ in size, era, variables, measurement practices, and patient composition.  Their proxy ontologies and DAGs are dataset specific, and nominally similar states are not necessarily equivalent.  Numerical estimates were therefore not pooled.  Operation of the workflow on both sources is positive evidence of engineering and methodological portability, especially given the inclusion of several predictive and causal model families.  It is not evidence that the estimates transport to other hospitals, countries, care practices, EHR systems, or time periods.

No prospective validation, independent external clinical validation, or formal transport analysis was performed.  Subgroup performance and effect stability by age, gender, ICU type, or other characteristics were not evaluated as a fairness study.  Consequently, the analysis does not establish clinical generalizability or subgroup equity.

\subsection{Reproducibility and Provenance}
\label{subsec:discussion:reproducibility-provenance}

The validated result manifest, source tables, source paths, and checksums provide strong numerical-transcription support for the values reported in Chapter~\ref{chap:results}.  They also improve artifact traceability by distinguishing canonical sources, duplicates, and excluded artifacts.  This is not the same as computational rerun reproducibility.  The \texttt{final-results/} archive is ignored or untracked in the main repository; exact numbered causal configurations are unavailable; result-generating source versions and copy histories are incomplete; and the nested causal repository contains local modifications.

Predictive split provenance, checkpoint-to-export mapping, and copied-output lineage remain incomplete.  Raw and processed datasets are external, exact processed-pickle hashes and producing commands are unavailable, and some paths are absolute to the producing machine.  Final software, environment, and hardware records are also incomplete.  Thus the available records support archived numerical transcription and materially improve artifact traceability, while a clean-checkout rerun is not currently demonstrated.  Scientific reproducibility---replication of findings under independently reconstructed data, definitions, and analysis---is weaker still and requires new evidence rather than additional formatting of the archive.

\subsection{LLM-Assisted Design Layer}
\label{subsec:discussion:llm-provenance}

The LLM-assisted design layer contributed structured proposals for proxy concepts, binary decision rules, missingness considerations, and dataset-specific DAGs.  Potential advantages are the systematic elicitation of domain concepts, scalable generation of candidate proxy definitions and rules, parallel dataset-specific proposals, explicit prompt artifacts that improve design provenance, and conversion of otherwise unstructured reasoning into inspectable source-coded artifacts.  These advantages concern organization, traceability, and design productivity; the thesis did not compare them with clinician-only construction or measure whether LLM assistance improved clinical or causal correctness.

The main design risks include hallucinated or unsupported clinical assertions; sensitivity to prompt wording, supplied context, model version, and hosted-model updates; unknown system, browsing, and decoding settings; incomplete reproducibility of hosted-model behavior; and possible dependence on hidden training data.  Project selection can introduce anchoring and confirmation bias, especially when plausible generated explanations are mistaken for independent evidence.  The prompt record does not establish complete accepted/rejected decisions, and only high-level rather than line-by-line prompt-output-to-code traceability is available.

Validation gaps remain substantial.  There is no documented clinician-panel adjudication of every proxy definition or graph edge, no chart-level proxy validation, no comparison of LLM-assisted with clinician-only or hybrid construction, no replication with another LLM family, and no prompt-perturbation or model-version sensitivity experiment.  Error can therefore propagate from an unsupported design proposal into a deterministic label, from that label into a prediction and aggregate, and from a misspecified exposure or graph into adjustment and effect estimation.  Preserved prompts and explicit source-code authority make the chain inspectable but do not validate it.

Predictive agreement with LLM-assisted rule-derived labels validates neither the clinical construct nor the LLM-generated design.  Sensitivity and permutation diagnostics do not validate the proxy ontology or DAG.  No claim is made that an LLM provides authoritative clinical judgment, discovered the true causal graph, or can replace a clinician or causal expert.

\subsection{Ethical and Clinical-Deployment Considerations}
\label{subsec:discussion:ethics-deployment}

Misclassification can misrepresent a record's clinical condition, amplify measurement bias, and generate misleading exposure contrasts.  Groups with different testing intensity or access to measurements may be affected disproportionately.  Automation bias is therefore a direct risk: users could overread model probabilities as diagnoses, proxy-state rankings as clinical priorities, adjusted estimates as treatment recommendations, or estimator agreement as causal certainty.

No dedicated fairness audit was performed.  The presence of age, gender, ICU type, or other variables in an adjustment set does not constitute fairness validation, and subgroup predictive performance and effect stability remain untested.  The repository does not establish the project-specific institutional approval, consent or waiver procedure, data-use agreement, or governance determination.  These facts must be supplied from authoritative institutional records and must not be inferred from dataset access, de-identification, or repository contents.

The deployment boundary is unambiguous: \emph{the system is a retrospective research workflow and is not suitable for patient-level clinical deployment or treatment recommendation in its current form}.  No prospective safety, workflow, impact, or human-factors evaluation was performed.  This limitation is not cured by predictive accuracy, estimator agreement, or artifact traceability.  Taken together, construct measurement, intervention definition, graph validity, temporal ordering, unmeasured confounding, support, uncertainty, transport, provenance, LLM-design reproducibility, fairness, and clinical validation remain the decisive threats to stronger interpretation.
